\documentclass[11pt]{article}

\usepackage{amsmath,amssymb}
\usepackage{xurl}
\usepackage[hidelinks]{hyperref}
\setlength{\emergencystretch}{2em}

\input{command}
\renewcommand{\ghissue}[1]{%
  \href{https://github.com/Dengnifer/MIPStarRE-A/issues/#1}{GitHub issue \##1}}

\title{Triangle Estimates and Isometry Transfer in Pauli Extraction}
\author{}
\date{2026-09-03}

\begin{document}
\maketitle

\section{At a glance}

\begin{itemize}
  \item \textbf{Difficulty.} Medium.  The extraction proof contains two
  incorrect numerical steps and later suppresses the range projection when
  passing from unitary conjugation to isometry conjugation.
  \item \textbf{Estimated weight.} A corrected positive-contraction estimate,
  a small-error case split for normalization, and one range-projection transfer
  lemma for a projective measurement family.
  \item \textbf{Mathlib/project split.} Positivity, contraction inequalities,
  isometry range projections, and norm estimates belong to finite-dimensional
  Hilbert-space algebra.  The state-dependent family distance and register
  placements are project-local.
  \item \textbf{Key mathematical inputs.} The inequality
  $H^2\geq2H-I$ for a positive contraction, the triangle inequality, and
  $Q=\phi\phi^\dagger$ for an isometry.  The repair is tracked in
  \ghissue{19}.
\end{itemize}

\section{Key theorem forms}

Let $H_X$ and $H_Z$ be the positive contractions of the extraction proof, and
let $\vartheta$ be a unit vector such that
\[
  \langle\vartheta,H_W\vartheta\rangle\geq1-\delta/2
  \qquad (W\in\{X,Z\}).
\]
Since $(I-H_W)^2\geq0$,
\[
  H_W^2\geq2H_W-I,
  \qquad
  \langle\vartheta,H_W^2\vartheta\rangle\geq1-\delta.
\]
Also $\lVert(I-H_W)\vartheta\rVert^2\leq\delta$, whence
\[
  \lVert H_X\vartheta-H_Z\vartheta\rVert^2\leq4\delta.
\]
Together with $H_XH_Z$ equal to the EPR projection, these estimates give the
required overlap bound, with constants absorbed into the stated
$O(\sqrt\delta)$ error.

For the final transfer, let $\phi\colon\mathcal H\to\mathcal K$ be an isometry
and set $Q=\phi\phi^\dagger$.  If $N_h$ is the unitary-conjugated extension of
a measurement operator $M_h$, then the exact identity is
\[
  \phi M_h\phi^\dagger=N_hQ=QN_h,
\]
not merely $N_h$.  Thus comparison of $N_h$ with an ideal projector must be
combined with control of $(I-Q)\Delta$ on the ideal comparison state
$\Delta$.

\section{Numerical defects in the source proof}

Lines~1715--1742 of
\path{references/qpbt-paper/14_analysis_of_the_pauli_basis_test.tex}
construct $H_X,H_Z$ and identify $H_XH_Z$ with the EPR projection.  The norm
estimate at lines~1743--1750 gives
$\lVert(I-H_W)\vartheta\rVert\leq\sqrt\delta$.  At lines~1753--1757 the
source substitutes $\delta$ inside the reverse triangle inequality
\cite[lines~1715--1783]{Ji2020MIPStar}.  That replacement is not valid for
small $\delta$.  The desired lower bound on
$\langle H_W^2\rangle$ follows directly from $H_W^2\geq2H_W-I$, as above.

The triangle inequality yields
\[
  \lVert H_X\vartheta-H_Z\vartheta\rVert\leq2\sqrt\delta,
\]
so its square is at most $4\delta$.  Line~1761 instead bounds that square by
$2\sqrt\delta$ without a stated restriction.  This inequality follows when
$\delta\leq1/4$; for larger error the theorem is made trivial by enlarging the
universal constant.  The same small-error split ensures that the projected
vector defined at lines~1769--1773 is nonzero before it is normalized.  These
are local numerical repairs and do not weaken the extraction conclusion.

\section{The range-projection obligation}

Lines~1827--1858 prove state-dependent closeness for the concrete operators
obtained by conjugating the original measurement with the swap unitaries.
Lines~1864--1875 then define isometries by adjoining EPR pairs and applying
those unitaries, and immediately state the corresponding closeness for
$\phi M_h\phi^\dagger$.  The omitted range projection makes this implication
non-definitional \cite[lines~1827--1875]{Ji2020MIPStar}.

Let $\Psi$ be the original state and suppose $\phi\Psi$ is close to the ideal
state $\Delta$.  Since $Q\phi\Psi=\phi\Psi$ and $Q$ is a contraction, the
triangle inequality bounds $(I-Q)\Delta$ by the state-extraction error.  Using
$\phi M_h\phi^\dagger=N_hQ$, split
\[
  (N_hQ-T_h)\Delta
  =(N_h-T_h)\Delta-N_h(I-Q)\Delta,
\]
where $T_h$ is the ideal Pauli projector.  Summing squared norms over $h$
uses projectivity and orthogonality to keep the second term under control.
The same argument is required for each player and each Pauli basis.

The finite-dimensional estimate is now formalized in
\path{MIPStarRE/QPBT/Test/Soundness/RangeProjection.lean}. For a complete POVM
$\{N_h\}$, an orthogonal projection $Q$, and $Q\vartheta=\vartheta$, it proves
\[
  \sum_h\|(N_hQ-T_h)\Delta\|^2
  \le 2\sum_h\|(N_h-T_h)\Delta\|^2+2\|\Delta-\vartheta\|^2.
\]
The proof uses $\sum_h N_h^\dagger N_h\le I$, so it applies to complete
POVMs as well as projective measurements. The same module proves the
isometry form under the exact intertwining identity $N_h\phi=\phi M_h$,
including the two local placements on an arbitrary bipartite state.
These estimates have no dependence on an assumed global polynomial-pair
measurement or on an extraction theorem.

The exact intertwining identity for adjoining a unit ancilla and applying a
unitary is proved in
\path{MIPStarRE/QPBT/Test/Soundness/Ancilla.lean}. In the extraction register
order, it constructs $J\psi=\psi\ot\eta$ and $\phi=UJ$, and proves
$U(M\ot I\ot I)U^\dagger\phi=\phi M$ from $\|\eta\|=1$ and
$U^\dagger U=I$. Thus this algebraic identity is available without a
separate intertwining assumption for these explicitly constructed maps.

Their application to supplied extraction data is now proved in
\path{MIPStarRE/QPBT/Test/Soundness/StateTransfer.lean} and
\path{MIPStarRE/QPBT/Test/Soundness/OperatorTransfer.lean}.
The local isometries adjoin the local EPR vectors and apply the concrete swap
unitaries. Reassociation identifies their joint action with the transformed
six-register state and identifies both ideal Pauli readouts exactly.
If the extraction data bounds the squared state error and both complete
unitary-conjugated Pauli comparisons by $\delta$, the resulting unsquared
state error is at most $\sqrt\delta$ and each isometry-conjugated Pauli
comparison has squared error at most $4\delta$. The actual transformed
state is fixed by both local range projections; no such condition is imposed
on the ideal state. The common bound is therefore
$\max\{\sqrt\delta,4\delta\}$, with no outcome or dimension factor.
These are conditional transfer results in a projective setting. Supplying
their extraction premise remains necessary before applying them to the
source-labelled soundness theorem.

\section{Returning through a general Naimark dilation}

The preceding ancilla argument must be distinguished from the dilation of an
arbitrary original POVM. Let $J$ be its ground embedding and $P_h$ its dilated
projectors. Naimark dilation gives $M_h=J^\dagger P_hJ$; it does not give
$P_hJ=JM_h$. If $\phi$ is a supplied soundness isometry for the dilated
strategy, set $Q=\phi JJ^\dagger\phi^\dagger$ and
$N_h=\phi P_h\phi^\dagger$. The original conjugated effect is exactly
$QN_hQ$. Summation over fibers shows that the same identity holds for the
actual Pauli readouts after the game's answer postprocessing.

Let $\theta$ be the transformed actual state and $\Delta$ the unit ideal
auxiliary-state/EPR vector. Exact dilation preparation preserves the state
error and gives $Q\theta=\theta$. The required decomposition is
\[
 (QN_hQ-T_h)\Delta
 =Q(N_h-T_h)\Delta-QN_h(I-Q)\Delta-(I-Q)T_h\Delta.
\]
On an EPR vector, a matrix on one register equals its transpose acting on
the other register. The opposite-register family therefore commutes with
$Q$ and controls the final term by $\|(I-Q)\Delta\|^2$ after summing squares.
The general complex-matrix transpose identity is proved explicitly. For the
existing fixed-field Pauli API, its previously proved
\path{pauliProj_transpose} identifies the transpose with the same Pauli
label. This uses the existing model, not an additional field restriction or
an assumption that arbitrary complex projectors are real.

If the state norm error and dilated Pauli squared-sum error are at most
$\delta$, the result is at most $3\delta+6\delta^2$. The families $N_h$
are generally not complete on the witness space: their total is a range
projection. Nevertheless isometry conjugation preserves
$\sum_h N_h^\dagger N_h\le I$. Thus any two such contractions have
squared distance at most four on a unit vector. The two estimates give the
common operator bound $9\delta$ for every $\delta\ge0$, while the state
norm error remains at most $\delta$. No ideal-state range assumption or
outcome-count factor occurs.

These calculations, both player placements, the actual value-preserving
dilation, and the pullback of supplied witnesses are proved in
\path{MIPStarRE/QPBT/Test/Soundness/NaimarkTransfer.lean}, with analytic and
EPR support in \path{CompressionBounds.lean} and \path{IdealPauli.lean}.
Replacing $a$ by $9a$ absorbs the factor nine into
$\delta_{\mathrm{qld}}^{a,b}$ for $a\ge1$ and $\eps\ge0$, preserving $b$.
The blueprint entries \path{ex:pauli-naimark-support} and
\path{ex:pauli-naimark-error-support} explicitly assume soundness data for
the dilated strategy. They do not construct those upstream data.

\section{Blueprint and formalization status}

The proof of Lemma~\path{lem:qld-unitary} and the soundness proof in
\path{blueprint/src/chapter/ch16_qpbt_extraction.tex} use the corrected
positive-contraction estimate, make the small-error case distinction explicit,
and isolate the range-projection calculation.  The chapter therefore preserves
the source's state and observable conclusions while exposing the omitted
argument.

The structure \path{ExtractionWitness} in
\path{MIPStarRE/QPBT/Extraction/Unitary.lean} contains only the two-sided
unitarity equations for the concrete swap maps, a normalized auxiliary state,
the transformed-state estimate, and the two-player Pauli-projector comparison.
The theorem \path{exists_extractionWitness_ofGlobalPairWitness} concludes the
existence of this data at the explicit extraction scale, but it takes the
\path{GlobalPairWitness} of Lemma~\path{lem:qld-4-7} as an explicit hypothesis
rather than obtaining it from \path{exists_globalPairWitness}.  That bundle
carries the projective global polynomial-pair measurements and both
point-consistency estimates, and the blueprint proof of
Lemma~\path{lem:qld-unitary} obtains it from Lemma~\path{lem:qld-4-7}; deleting
the hypothesis would leave that construction unproved.  By the criterion of
\path{docs/paper-gaps/proof-gap-protocol.tex} the hypothesis is therefore proof
debt rather than a faithful boundary condition, and the declaration is a
conditional helper: it does not certify the source lemma, and the blueprint
node Lemma~\path{lem:qld-unitary} accordingly neither links it nor carries any
statement-level formalization mark.  The declaration is recorded instead on the
separate formalization-support node
\path{lem:qld-unitary-given-global-pair}, whose statement assumes the global
polynomial-pair measurement explicitly and whose formal link points only to the
conditional helper.  It therefore cannot be mistaken for the source result.
Apart from that hypothesis the
conclusion matches the Stage~4.2 extraction lemma, including the corrected
numerical estimates. The general range-projection estimates described above
justify the analytic part of the later passage to isometries. The concrete
isometries, intertwining, and state and measurement transport from supplied
extraction data are proved by the two transfer modules above. The separate
blueprint support entries
\path{ex:pauli-supplied-extraction-support} and
\path{ex:pauli-supplied-extraction-error-support} explicitly assume this
data in a projective setting. The latter composes the extraction error
estimate with the square-root and prefactor adjustments for all three
conclusions. Neither entry certifies the unrestricted source theorem.

The conditional character of the result is visible wherever the declaration is
invoked: the suffix \path{ofGlobalPairWitness} names the extra premise in the
declaration itself, and its docstring records the statement as conditional and
its dependence on the assumed witness as unfaithful.  Discharging the debt
requires composing \path{exists_globalPairWitness} with this helper to obtain
the source lemma; no declaration carries out that composition.

Issue~\#19 owns this concrete Stage~4.2 interface and discrepancy record.
\ghissue{123} owns the source-facing composition through
\path{exists_globalPairWitness}.  \ghissue{47} tracks the separate
range-projection transfer needed by \path{pauli_soundness}; that theorem must
be derived from the concrete data rather than stored in the witness.
The independent range-projection estimates and the complete transfer from
supplied extraction data are recorded by \ghissue{529}; these results do not
discharge the global-pair construction, the zero-error endpoint of that
construction, or the source-domain and combined-lines error comparisons.
The separate return through the genuine Naimark dilation is now proved as
above, conditional on supplied soundness data for that dilation.

\section{Conclusion}

The printed triangle calculation has repairable numerical errors, and the
last passage from unitaries to isometries omits a genuine range-projection
argument.  Neither defect is a counterexample to Pauli soundness.  The former
is corrected by positive-contraction inequalities and a small-error split;
the latter has a proved general analytic estimate and a proved application to
supplied concrete extraction data. Independently of both defects, the
present extraction declaration is conditional on the global polynomial-pair
measurement, and that conditionality must be removed before it may be
advertised as the source lemma.

\bibliographystyle{alpha}
\bibliography{references}

\end{document}
